% Named-mode design, aligned with docs/design_spec.md §2.
% Named read-only access; the common object representation and lifetimes are in Section 3.
% Historical results and provisional fixed-pool data do not evaluate the CPU-local design.
\section{Named Mode: Objects Without a Mapping}
\label{sec:named}

Named mode provides read-only handle operations without mappings. It uses the shared objects
and lifetimes of \S\ref{sec:lending:objects}.

\subsection{Interface}
\label{sec:named:interface}

An application registers a \emph{store} on an \iouring ring, specifying its retention limit,
largest object size, and maximum numbers of objects and streams. It attaches connected TCP
sockets to the store. Each attachment returns a \emph{stream token} and gives the store
exclusive receive access to that socket. Operations identify stream bytes by absolute offset,
starting at the first unread byte when attachment occurs.

A parser first inspects a header, then discards the header bytes and names the body range.
\code{RECV\_OBJECT} can be posted before the body arrives, but returns a handle only when the
whole range is present. The main operations are:

\begin{center}
\footnotesize
\setlength{\tabcolsep}{3pt}
\begin{tabular}{@{}
  >{\raggedright\arraybackslash}p{0.32\columnwidth}
  >{\raggedright\arraybackslash}p{0.64\columnwidth}@{}}
\toprule
Operation & Effect \\
\midrule
\code{INSPECT} & Copy a bounded window of stream bytes without consuming it. \\
\code{AVAILABLE} & Count contiguous available bytes without copying or consuming them. \\
\code{RECV\_OBJECT} & Take an exact stream range as a complete object. \\
\code{DISCARD} & Consume an unwanted stream range. \\
\code{SEND} & Send object ranges, optionally with generated framing. \\
\code{FREE} & Remove an object handle. \\
\code{READ\_OBJECT} & Copy an exact, bounded range of a completed object. \\
\code{CREATE} & Copy application bytes into a new object. \\
\bottomrule
\end{tabular}
\end{center}

\code{INSPECT} may return only part of a header. The parser accumulates it by requesting the
first missing byte. A failed \code{RECV\_OBJECT} publishes no partial object and returns its
captured prefix to the stream while the stream remains open. Discard is different: a failed
or canceled \code{DISCARD} may already have consumed bytes, which cannot be restored.

The store, rather than the source socket, holds completed objects, so they survive stream
closure. \code{STREAM\_CLOSE} releases the stream token, undecided input and receive claim;
closing the application's socket descriptor alone does not close that stream. A store can
be exported through a file descriptor and imported by other rings. They share its handles,
capacity and per-socket send ordering.

\subsection{Capturing received bytes}
\label{sec:named:capture}

The \emph{collector} removes bytes from TCP's receive queue only as it retains or discards
them. Stable, readable page fragments are retained by reference. When it consumes an unshared
packet completely, the collector can transfer the packet's page references rather than
incrementing their counts. Where this transfer is not possible, it takes its own references.

Data that cannot safely be retained is copied into kernel-owned buffers. This includes
payload in non-page allocations and shared or cloned buffers whose contents another owner
could change, as can occur over loopback and veth. Unreadable fragments, such as device-memory
payload, are refused. Capture counts fallback copies; retaining ordinary TCP data does not
guarantee that every byte is taken by reference.

Bytes not yet assigned to an object or discarded are \emph{undecided}. They occupy store
entries while the application parses. Collection follows demand: it can take a batch when no
undecided bytes remain, continue for a waiting inspection, or advance toward a pending range
decision. It stops unnecessary lookahead once undecided input awaits a decision. Unconsumed
bytes stay in TCP, so a slow parser eventually closes the receive window. Exhausting store
entries also stops capture (\S\ref{sec:named:memory}).

\subsection{Sending object ranges}
\label{sec:named:send}

A send walks the object's segments and passes page references to TCP through the path used
by \splice and \sendfile. If the socket fills, the request keeps its position and resumes at
the first unqueued byte. One request can include a generated prefix and suffix, which are
copied, or gather ranges from several objects in a vectored send. A queue for each store and
destination socket keeps these frames in order across retries, including sends from rings
sharing the store. Other writers do not participate in that queue.

A send completion reports bytes queued to TCP, including framing. It does not report remote
delivery or page reclamation, and there is no later buffer-reuse notification for the object
body. A send marked as the object's last use removes the handle at admission. Its completion
reports whether the handle was consumed independently of whether the send succeeded. An
error or short send after that transfer does not restore the name. A frame that terminates
short also cancels its current followers; the application must recover the connection rather than
assuming that frame boundaries remain intact.

Two choices reduce the work per send. Vectored sends share request and socket-lock overhead
across ranges. A separate optimization to the stack's splice helper reduces unnecessary
splitting of eligible large-page ranges; it also benefits other users of that helper
(\S\ref{sec:impl}). Received objects can still contain many small fragments. The transmit
fragment limit may then require extra packets compared with a dense copying sender. Larger
packets, a higher fragment limit, or a dense copy trade different costs against that limit.
\todom{compare the received-fragment send configurations and their packet counts (Q6)}

\subsection{Memory limits and a full store}
\label{sec:named:memory}

Each segment holds at most one receive-buffer reference, so $N$ segments retain at most $N$
buffers across undecided input, pending objects and completed objects. A buffer split
between objects consumes separate segments, so capacity can be exhausted before that many
distinct buffers are held. The head array independently limits the number of pending and
completed objects.

The base segment count is the configured byte limit (\code{hard\_limit}) divided by the
system page size. Each pool adds enough entries to cover all CPU-local stock: four batch
capacities per possible CPU. This prevents stock in other caches from causing allocation
failure before the configured count is in use. Live usage can exceed the base count by at
most that extra stock; the physical arrays, not \code{hard\_limit} alone, bound retention.

A short slice retains its whole backing buffer, so the limit is not useful payload capacity.
Larger driver buffers also increase the physical footprint per segment. We call retained
space outside useful payload \emph{slack}.

Registration reserves the configured retention limit and charges it and the provisioned
arrays to the registering task's memory control group once. This retention reservation is
not an allocation of payload pages. Owned-copy buffers require a separate pool, including
its extra CPU-local stock, charged and reserved at setup (\S\ref{sec:impl}). Taking an
entry adds no per-operation charge. Importing a store adds neither capacity nor a new owner.
These limits exclude rings and socket memory; TCP can retain pages after the store releases
them.

\paragraph{When capacity runs out.} Requests do not wait for a free head or segment. A request
that can obtain one from neither its local pair nor the shared pool fails at once with \code{-ENOMEM}.
Free entries in other pairs are not reclaimed before teardown. The collector stops, keeps bytes already
captured and reports a nonterminal \code{-ENOMEM} event; bytes it has not taken stay in TCP.
A \code{RECV\_OBJECT} that cannot finish returns its captured prefix to undecided input and
publishes no handle.

The application must free objects or discard unwanted input, then retry. Its next stream
request runs another collection pass. \code{FREE} does not wake a queue of memory waiters,
and a pending discard can advance without capturing more payload. An application that only
waits for an object's completion while ignoring the collector's memory event can stall
itself. Eviction and the number of outstanding requests remain application policies.

\paragraph{Slack remains a tradeoff.} The arrays do not grow, and this design does not yet
compact held data. Smaller receive buffers and better packet coalescing can reduce slack.
Compaction would copy retained payload into denser buffers at an explicit copy cost. Earlier
builds support it; its cost needs a separate comparison (\S\ref{sec:eval:fwd}).

\subsection{Safety and shared authority}
\label{sec:named:safety}

The kernel checks a handle's store, generation and byte bounds. Generation checks prevent a
stale handle from naming a new object in a recycled entry.
Exclusive receive access prevents competing readers, sockmap programs or kernel TLS from
changing the stream's byte offsets. Authorization to read the socket is checked at
attachment.

The named mode needs no administrative privilege, but remains subject to registration and
socket permissions. Exporting a store deliberately shares authority: an importing ring can
read, send or free its objects within the same capacity. It does not receive an isolated
namespace. Closing streams, canceling requests and tearing down stores release their
references without invalidating references still held by admitted sends or TCP. Finally,
immutability does not protect against a trusted process that maps the device's receive
memory (\S\ref{sec:mapped}).
